% Generated draft: requires mathematical and source review.
\documentclass[11pt]{article}
\usepackage[margin=25mm]{geometry}
\usepackage{fontspec}
\setmainfont{Latin Modern Roman}
\usepackage{amsmath,amssymb,mathtools,unicode-math}
\setmathfont{Latin Modern Math}
\usepackage{enumitem,needspace,xurl,hyperref}
\setlength{\parindent}{0pt}
\setlength{\parskip}{5pt}
\setlength{\emergencystretch}{3em}
\newcommand{\N}{\mathbb{N}}
\newcommand{\Z}{\mathbb{Z}}
\newcommand{\Q}{\mathbb{Q}}
\newcommand{\R}{\mathbb{R}}
\newcommand{\C}{\mathbb{C}}
\newcommand{\F}{\mathbb{F}}
\newcommand{\A}{\mathbb{A}}
\newcommand{\PP}{\mathbb{P}}
\newcommand{\E}{\mathbb{E}}
\newcommand{\eps}{\varepsilon}
\newcommand{\dd}{\,\mathrm{d}}
\begin{document}
\section*{UnsolvedMath ID: 700653}
\textbf{Source problem code:} catalogue entry
\subsection{Definitions and mathematical statement}
For each integer \(n\geq 1\), let \(R_n=\Z[x_1,\ldots,x_n]\) be the polynomial ring in \(n\) commuting indeterminates, and put \(m=n+3\). Let \(I_m\) denote the identity matrix and let \(e_{ij}\) denote the \(m\times m\) matrix unit. Define
\[
T_n=\{I_m+a e_{ij}:a\in R_n,\ 1\leq i,j\leq m,\ i\neq j\},
\qquad E_m(R_n)=\langle T_n\rangle\leq\operatorname{SL}_m(R_n).
\]
Here the elements of \(T_n\) are the elementary transvections; their inverses also belong to \(T_n\). A group has finite breadth with respect to this generating set if there is a finite bound on the number of such factors needed to express any of its elements.

Under this interpretation of the source's terminology, the question is whether
\[
\forall n\in\Z_{\geq 1}\ \exists B_n\in\Z_{\geq 0}\ \forall g\in E_{n+3}(R_n)\ \exists \ell\in\{0,\ldots,B_n\}\ \exists t_1,\ldots,t_\ell\in T_n:
\qquad g=t_1\cdots t_\ell.
\]
The empty product is the identity. The bound may depend on \(n\), but must be independent of \(g\); no bound uniform in \(n\) is requested by this formulation.

The supplied record does not define ``transvections'' or ``finite breadth.'' The displayed formulation uses elementary transvections and bounded product length. Whether a broader class of transvections, such as conjugates of elementary transvections, was intended requires source review.
\subsection{Short English statement}
For each integer \(n\geq 1\), can every element of the elementary matrix group over \(\Z[x_1,\ldots,x_n]\), in dimension \(n+3\), be written as a product of a bounded number of elementary transvections? The bound may depend on \(n\). This is a statement-only interpretation of the supplied question; its terminology requires human review.
\subsection{Sources}
\begin{itemize}
\item Supplied local UnsolvedMath registry record, permanent ID \(700653\), exact source problem code \texttt{catalogue entry}, titled ``Kourovka Notebook Problem 10.38.'' The record was created and updated on 2026-10-16 and labels the problem \texttt{open}; that label is supplied metadata, not an independently verified mathematical or literature assessment.
\item The supplied background identifies \emph{The Kourovka Notebook (Unsolved Problems in Group Theory)}, 10th issue (1986), problem 10.38, and the 21st issue edition (2026), identified as \texttt{arXiv:1401.0300v48}. The registry describes the collection as edited by E. I. Khukhro and V. D. Mazurov. These source identifications are supplied and have not been independently checked.
\item The background and proposer field attribute the problem to G. A. Noskov, with proposed year 1986. The statement separately begins ``(W. van der Kallen).'' Both attributions are retained; their precise relationship is not explained by the record.
\end{itemize}
\end{document}
